\begin{abstract}
Low-rank adaptation can alter the safety behavior of an aligned language model, but the mechanism that makes such behavior persist or rebound during continued fine-tuning remains unclear. We study this problem as an empirical mechanism question rather than proposing a new safety method. In a controlled \modelname{} risk-LoRA setting, unsafe adapters exhibit reproducible spectral and output-side subspace structure, yet global spectral quantities, post-hoc geometric interventions, training-time spectral penalties, and tested local scalar proxies do not recover safe behavior. Sequence-level safety supervision is behaviorally active across three early-training seeds. In a later seed-1 comparison, maintaining a 20\% natural safety replay rate reduces final unsafe behavior from 0.56 to 0.28 relative to a schedule that decays to 10\%, but the stronger persistence control reaches only 0.72 prompt-relevant redirect coverage and exhibits a 0.72 mechanical-template share under the historical audit. These observations motivate a stronger identification test that holds cumulative supervision dose fixed while varying its temporal placement, followed by frozen prompt-aware quality measurement, independent human validation, a second setting, and one prompt-disjoint confirmation. \PaperPending{replace this design sentence with the registered three-seed equal-dose result only after all promotion gates pass.} The current evidence supports a distinction between diagnostic geometry and behavioral control, and identifies persistent sequence supervision as a strong control with a nontrivial safe-response-quality bottleneck; it does not validate a new training algorithm.
\end{abstract}
